\section{Introduction}
\label{sec:introduction}

Homologous recombination, which we will hereafter refer to simply as recombination, substantially contributes to the diversification of most known viruses and cellular lifeforms. 
The exact primary reason(s) that recombination is so rampant in the biosphere are unsettled but are likely a combination of: (1) it is a core process in the repair of double stranded DNA breaks in cellular lifeforms, (2) it enables the "rewinding" of maladaptive evolution in that it can help purge genomes of harmful mutations, and (3) it enables the "fast-forwarding" of evolution in that it can concentrate distinct adaptive mutations within individual genomes \citep{simonloriere2011recombination,didelot2010impact}. 
The fact that biological processes have evolved to foster the occurrence of homologous recombination in a broad array of different RNA viruses where it plays no role in the repair of broken genome strands, suggests that, at least in some instances, its primary adaptive value is in the shuffling of mutations between one generation and the next. Nevertheless recombination rates vary widely between different viruses \citep{simonloriere2011recombination} and cellular lifeforms \citep{didelot2010impact,mcvean2004fine} and, when it does occur, it tends to occur more frequently in some genome regions than in others \citep{mcvean2004fine,lytras2022exploring}. Further, contingent on the precise sites in a genome where it occurs, it has a higher or lower tendency to yield recombinant offspring genomes that are either fitter or less fit than their parental genomes \citep{simonloriere2011recombination,lytras2022exploring}.

In the context of pathogenic microbe evolution, these recombinant offspring genomes could have altered transmission dynamics, cellular tropisms, host ranges, pathogenicity, antigenicity, immune-evasiveness and antimicrobial resistance \citep{alfonsi2024data,didelot2010impact,croucher2011rapid}. In sexually reproducing eukaryotes, across successive generations highly orchestrated inter-chromosomal recombination during meiosis enables the purging of deleterious mutational load and the assembly of hybrid genomes with complex novel phenotypes \citep{rieseberg2003major,mcvean2004fine}.

While interesting in its own right, recombination is also a nuisance in the context of phylogenetics-based molecular evolution analyses. Specifically, it undermines the core assumption of phylogenetics: that the evolutionary relatedness of a set of homologous sequences can be accurately described using a bifurcating tree. If recombinant genomes have tracts of sequence inherited from parental genomes that had sufficiently distinct evolutionary histories, attempts to place the recombinant within a single bifurcating phylogeny will tend to shorten deep inner branches, lengthen outer branches and cause spurious branching patterns \citep{schierup2000consequences,kosakovsky2006gard}.
These recombination-induced tree distortions will at least mildly impact the downstream inference of molecular clock rates, the estimation of node dates, ancestral sequence and ancestral character state inference, transmission cluster reconstructions, and the inference of variations in historical population sizes \citep{schierup2000consequences,posada2002evaluation}. When using phylogenetic codon models to analyse patterns of selection based on rates of synonymous and nonsynonymous substitutions, the apparent homoplasies caused by recombination can mimic convergent amino acid substitutions across multiple branches of a phylogenetic tree.
Under standard selection detection tests, including MEME and BUSTED \citep{murrell2012detecting,murrell2015gene}, recombination produces spurious signals of diversifying positive selection \citep{kosakovsky2006gard}.

Therefore, detecting recombination breakpoint sites and identifying recombinant sequences within nucleotide sequence alignments has become a routine task in phylogenetics-based molecular evolution analyses. Identification of phylogenetically disruptive recombination breakpoint sites at the scale of whole alignments with recombination detection methods such as GARD \citep{kosakovsky2006gard} can be used to partition alignments into recombination-depleted segments that can then be used individually for downstream phylogenetics-based evolutionary inference. If phylogenetically disruptive recombination breakpoints are so pervasive that recombination-depleted segments are so small that they cannot be used for accurate phylogenetic inference, "triplet scanning" recombination detection methods like 3SEQ \citep{lam20183seq} or many of those implemented in RDP5 \citep{martin2021rdp5} can be used to individually scan every combination of three sequences in an alignment for evidence of recombination and identify breakpoints within individual sequences. With both identified recombinant sequences and breakpoint locations in hand, individual sequences can then be segmented or removed to yield full-length recombination-depleted alignments for downstream analyses.  

These two different approaches to accounting for and managing the confounding effects of recombination on molecular evolutionary analyses foreshadow a much deeper divide between the ways in which the different detection methods work. 
GARD is an example of a phylogenetic change-point method that formulates recombination breakpoint detection as a model selection task across partitioned trees under continuous-time Markov substitution models. GARD finds optimal breakpoint boundaries through likelihood ratio criteria while fitting independent branch lengths across candidate alignment segments. Because the joint combinatorial space of partition boundaries and tree topologies scales exponentially with breakpoint count and cubically with taxon count ($\mathcal{O}(L^K \cdot N^3)$ for $K$ breakpoints across $N$ taxa of alignment length $L$), multi-breakpoint searches across hundreds of genomes with GARD require days of cluster compute.

Triplet scanning detection methods like those implemented in 3SEQ and RDP5, forgo computationally expensive evolutionary modeling. Instead they rely on the heuristic detection of localized relative or absolute differences in the similarities of sequence pairs within the context of individual sequence triplets drawn from a bigger alignment. While evaluating rapidly across alignments of a few hundred sequences, exhaustive triplet scanning approaches scale poorly $\mathcal{O}(N^3 L)$ to the point that they become untenable when alignment sizes increase to over a few thousand sequences.
Since the individual sequence triplets that display evidence of recombination have no phylogenetic reference frame against which to decide which sequence in the triplet is the recombinant (and which are parent-like), triplet scanners must be augmented with downstream phylogenetic methods to make these calls; which can become very difficult to make accurately when close relatives of both of a recombinant sequence's parental lineages are not present in an analyzed alignment. Downstream phylogenetic methods and other techniques (such as iteratively deconstructing detected recombinant sequences into their parental components and re-scanning for recombination) are also needed to determine the counts of actual ancestral recombination events that underlie detected recombination signals; because any given ancestral recombination event might be independently detectable with tens or hundreds of the different sequence triplets drawn from an analyzed alignment. Finally, if triplet scanning methods do not take steps to phylogenetically contextualize detected recombination signals they run the risk of identifying lineage-specific increases in substitution rates within defined genome sub-regions as evidence of recombination. All of this additional phylogenetic contextualization incurs extra computational costs over-and-above the triplet-scanning recombination signal detection components of 3SEQ and RDP5.

Here we introduce \RhizAeon, an extremely fast and linearly scaling open-source recombination detection tool that retains phylogenetic context without the construction of discrete phylogenetic trees. Its name derives from the Greek \textit{rhiza} (root) and evokes the interconnected network of the biological rhizome across deep evolutionary time (\textit{aeon}) to supplant strictly bifurcating arboreal trees.


